\documentclass[12pt,a4paper]{article}
\usepackage{graphicx}
\usepackage{titlesec}
\usepackage{setspace}
\usepackage{geometry}
\geometry{margin=1in}

\usepackage{graphicx}
\usepackage{subcaption}
\usepackage{float}
\usepackage{hyperref}

% Formatting
\setstretch{1.15}
\titleformat{\section}{\large\bfseries}{\thesection.}{1em}{}
\titleformat{\subsection}{\normalsize\bfseries}{\thesubsection.}{1em}{}
\hypersetup{colorlinks=true, linkcolor=blue, citecolor=blue, urlcolor=blue}


%----------------------------------------
\begin{document}

% --------- First Page (Cover Page) ---------
\begin{titlepage}
    \centering
    \includegraphics[width=3cm]{Logo.jpeg}\par\vspace{1cm}
    {\LARGE \textbf{Department of Electrical and Information Engineering}}\\[0.4cm]
    {\Large \textbf{University of Ruhuna}}\\[1.5cm]
    {\Large \textbf{Group Project Assignment}}\\[0.2cm]
    {\large \textbf{Advanced Artificial Intelligence}}\\[0.2cm]
    {\large \textbf{EC7203}}\\[6cm]

    % Centered student and project info
    \begin{tabular}{rr}
        \large \textbf{EG/2021/4384} & \large \textbf{Achintha H.G.R.} \\
        \large \textbf{EG/2021/4511} & \large \textbf{Fernando S.A.U.} \\
        \large \textbf{EG/2021/4516} & \large \textbf{Fonseka C. M.} \\
        \large \textbf{EG/2021/4751} & \large \textbf{Rathnaweera P. A.G. N.} \\
    \end{tabular}\\[1cm]
    {\large \textbf{Computer Engineering (CE)}}\\

    {\large \textbf{xx/05/2026}}\\[2cm]
    
\end{titlepage}


% --------- Table of Contents ---------

\clearpage
\phantomsection
\addcontentsline{toc}{section}{Contents}
\label{contents}
\tableofcontents
\clearpage

% --------- Sections ---------
%========================= INTRODUCTION ========================%
\section{Introduction}
 
Social media has become the dominant channel through which information --
and misinformation -- reaches the public, and the volume of unverified
claims now far exceeds what human fact-checkers can manually review
\cite{vykopal, factgpt}. This pressure has driven a
shift toward automated fact-checking, where large language models (LLMs)
extract claims, retrieve supporting evidence, and reason toward a verdict
\cite{veracity, bangare}. A parallel line of work
decomposes this process across cooperating or debating agents, improving
transparency and reducing the brittleness of a single monolithic model
\cite{gautam, ted, mace}.
 
These advances, however, are overwhelmingly built for English and other
high-resource languages. Recent evidence shows that LLM performance on
fact-checking tasks degrades sharply for low-resource languages and
non-Latin scripts, and that the common workaround -- translating everything
into English -- introduces its own loss of meaning
\cite{vykopal}. The problem is acute in Sri Lanka, whose
information ecosystem is genuinely trilingual (Sinhala, Tamil, and English)
and where social media text is routinely \emph{code-mixed}: a single post
may switch scripts mid-sentence or render Sinhala phonetically in Latin
characters (``Singlish''). Off-the-shelf, cloud-centric fact-checkers treat
such input as noise, and existing multilingual systems seldom address
code-mixing as a first-class concern \cite{hemtfake}.
 
This work presents \textbf{TruthBomb}, a local-first, multi-agent
fact-checking engine designed for Sri Lankan social media content. A
fine-tuned, deterministic extractor first \emph{normalizes} a noisy
code-mixed post into a single declarative English claim; a hybrid retriever
(dense + lexical) gathers evidence from a continuously ingested local
knowledge base; a deterministic, confidence-based router escalates to web
search only when local evidence is weak; an adjudication stage produces a
verdict; and a final stage renders that verdict back into English, Sinhala,
and Tamil. Rather than claiming broad superiority, this project makes three
\emph{targeted} contributions appropriate to its scope: (i) a dedicated
fine-tuned stage for claim extraction from code-mixed, low-resource input;
(ii) local-first OSINT retrieval with score-based web escalation; and
(iii) trilingual verdict generation for genuinely low-resource languages.
%==================================================================================%


%========================= MOTIVATION ===============================%
\section{Motivation}
\label{sec:motivation}
 
The case for a Sri Lanka--specific system rests on three observations that
generic fact-checkers do not accommodate.
 
\paragraph{Language reality, not language coverage.}
Many ``multilingual'' systems advertise support for hundreds of languages by
relying on machine translation into English before any reasoning occurs
\cite{bangare}. This pipeline assumes clean, monolingual
input. Sri Lankan posts violate that assumption: they mix Sinhala, Tamil,
English, and romanized Singlish within a single utterance, where a naive
translation step either fails or silently distorts the claim. We therefore
treat the first stage not as translation but as \emph{normalization} --
language conversion, claim isolation, and declarative reformulation in one
fine-tuned step -- so that downstream agents always receive a clean,
verifiable target. Because the same input distribution is what general LLMs
handle worst \cite{vykopal}, a small specialized model
fine-tuned on this exact task is both the cheapest and the most reliable
remedy \cite{factgpt}.
 
\paragraph{Evidence sovereignty, latency, and auditability.}
Cloud-only fact-checkers send every query to an external search or LLM
service, which is costly, slow, and opaque about \emph{which} evidence drove
a verdict. TruthBomb instead maintains a continuously ingested local
knowledge base and consults it first; web search is invoked only as an
escalation path when local confidence is low. Crucially, this routing is
\emph{deterministic and score-based} rather than delegated to an LLM, which
avoids an extra inference call and makes the decision auditable -- a property
prior work identifies as essential for trustworthy fact-checking
\cite{veracity}.
 
\paragraph{Accessibility for the people most exposed.}
A verdict is only useful if its intended audience can read it. The
communities most exposed to Sinhala and Tamil misinformation are precisely
those least served by English-only output. Producing the final verdict and
its reasoning in all three languages is therefore not a cosmetic feature but
the point of the system: it closes the loop between the languages
misinformation arrives in and the languages a correction must be delivered
in. In keeping with the project's scope, we frame this as a best-effort,
honestly evaluated capability rather than a claim of production-grade
translation quality.
%==================================================================================%


%========================= PROBLEM STATEMENT ========================%
\section{Problem Statement}
\label{sec:problem-statement}
 
\textbf{Core problem.} Given a short, potentially noisy and code-mixed social
media post written in Sinhala, Tamil, English, or romanized Singlish, produce
a verifiable, evidence-grounded verdict together with reasoning that is
accessible to the language community in which the post originated.
 
A monolithic LLM call is ill-suited to this task: it conflates
language understanding, claim isolation, evidence gathering, and judgment
into one opaque step that, as noted above, performs worst on exactly the
low-resource, code-mixed inputs we care about. We therefore decompose the
problem into four sub-problems, which correspond directly to the stages of
the proposed pipeline:
 
\begin{itemize}
  \item \textbf{P1 -- Claim extraction and normalization.} Convert a
  code-mixed, possibly non-declarative post into a single normalized English
  declarative claim, or emit a \texttt{NO\_VERIFIABLE\_CLAIM} sentinel when
  the post contains no checkable assertion. This must be deterministic and
  structurally constrained so downstream stages receive reliable input.
 
  \item \textbf{P2 -- Evidence retrieval and routing.} Retrieve relevant
  evidence from a local knowledge base using hybrid (dense + lexical)
  search with re-ranking, and decide -- via a deterministic,
  confidence-based rule -- whether local evidence is sufficient or whether
  the query must escalate to web search.
 
  \item \textbf{P3 -- Adjudication.} Combine the claim and the retrieved
  evidence into a verdict (e.g.\ supported / refuted / unverifiable) with an
  associated confidence and an audit trail of the evidence used, resolving
  ambiguous cases with a reasoning step only when rules are insufficient.
 
  \item \textbf{P4 -- Trilingual verdict generation.} Render the verdict and
  its reasoning in English, Sinhala, and Tamil so the output is accessible to
  the originating community.
\end{itemize}
%==================================================================================%
%========================= LITERATURE REVIEW ========================%
\section{Literature Review}
The detection of misinformation on social media has converged on a small set of
recurring design decisions: whether to \textit{classify} content or to
\textit{reason} over retrieved evidence; whether to use a single model or to
decompose the task across cooperating agents; how far to push genuinely
multilingual coverage; and how much of the verification process to expose to the
user. This section surveys eight recent works against these axes, identifies the
paradigm best suited to the present setting --- fact-checking code-mixed Sri
Lankan social-media content across Singlish, Sinhala, Tamil, and English --- and
positions the contribution of this project against the resulting gaps.

\subsection{The Research Landscape}

\subsubsection{Multi-agent verification pipelines}
The clearest recent trend is the decomposition of verification into specialised,
cooperating agents. Saha et al.~\cite{mace} verify claims against tabular data
using three agents: a Planner that devises a reasoning strategy, an Executor that
performs the verification, and a Verifier that audits the Executor and can return
a ``revise'' signal. Notably, this MACE framework relies on zero-shot
Chain-of-Thought prompting rather than fine-tuning, yet reaches state-of-the-art
results while using far smaller models (27--92B versus 235B parameters). Its
central lesson is that careful orchestration and an explicit verify-and-revise
loop can substitute for expensive specialisation, while producing interpretable
verification traces as a by-product.

Gautam~\cite{gautam} widens the remit, arguing that most agentic systems stop at
detection. The proposed framework employs five agents --- Indexer, Classifier,
Extractor, Corrector, and Verification --- that together cover classification,
evidence retrieval and ranking, correction, and, distinctively, source provenance
and root-cause identification. The work is conceptual, presenting the framework
without experimental evaluation, but it elevates audit trails and source
provenance to first-class design concerns.

Liu et al.~\cite{ted} take decomposition in an adversarial direction. In their
TruEDebate (TED) system, two ``DebateFlow'' teams argue for and against a claim
through opening statements, cross-examination, rebuttal, and closing, while an
``InsightFlow'' layer (a Synthesis agent plus an Analysis agent built on a
role-aware encoder and a debate graph) renders the final judgment. On the ARG-EN
and ARG-CN datasets, TED outperforms traditional methods while substantially
improving interpretability, demonstrating that structured disagreement --- not
just sequential delegation --- is a viable mechanism for resolving ambiguous
claims.

\subsubsection{End-to-end LLM fact-checking systems}
A parallel line of work prioritises deployable systems over novel agent
topologies. Curtis et al.~\cite{veracity} pair an LLM with web-retrieval agents in
their Veracity system to assess user-submitted claims, returning a numerical
veracity score alongside plain-language explanations, with multilingual support
and a messaging-style interface. Its emphasis on transparency, accessibility, and
serving as an open research test-bed reflects a recognition that high accuracy
alone does not make a system trustworthy or usable.

Bangare et al.~\cite{bangare} present a more applied effort built around the
Gemini API for claim understanding, translation, and reasoning, combined with the
Google Fact Check Tools API for verification, with LLM reasoning as a fallback
when no prior fact-check exists. Although its evaluation is limited, it is a useful
reference because it instantiates the same engineering pattern adopted here: claim
extraction, then LLM reasoning, then verification against external sources, across
multiple languages.

\subsubsection{Claim matching and previously-fact-checked claim detection}
Two works isolate the upstream problem of recognising claims rather than
adjudicating them. Choi and Ferrara~\cite{factgpt} treat claim matching as a
textual-entailment task --- deciding whether new content aligns with, contradicts,
or is irrelevant to an already-debunked claim --- and fine-tune smaller LLMs on a
synthetic dataset to match the accuracy of much larger models. Their FACT-GPT work
supplies two findings directly relevant to this project: that small, task-specific
fine-tuned models can outperform zero-shot large models, and that synthetic data
is a legitimate training resource, subject to the usual distillation caveats.

Vykopal et al.~\cite{vykopal} provide the most rigorous multilingual evidence in
the set: an evaluation of seven LLMs across twenty languages in monolingual and
cross-lingual settings. LLMs perform well on high-resource languages but degrade
on low-resource ones, and translating source text into English measurably helps
the low-resource cases. This empirically supports normalising low-resource and
code-mixed input toward English as an intermediate representation.

\subsubsection{Multimodal and multilingual classification}
Jadhav et al.~\cite{hemtfake} represent the alternative paradigm: rather than
reasoning over evidence, they classify. Their HEMT-Fake work contributes a
74{,}000-article multilingual--multimodal dataset (Hindi, Gujarati, Marathi,
Telugu, and English with paired images) and a hybrid architecture fusing
transformer embeddings, a CNN--BiLSTM encoder, ResNet image features, and
GraphSAGE metadata via multi-head attention, with an explainability layer
combining attention, SHAP, and LIME. It reports roughly a 5\% Macro-F1 gain over
XLM-R and mBERT (7--8\% on low-resource languages) and resilience to adversarial
paraphrasing. It shows that the classification paradigm can deliver strong
multilingual and multimodal performance with explanation, but it does not produce
evidence-grounded verdicts or reasoning traces.

\subsection{Selecting the Most Suitable Paradigm}
For the present setting, the agentic, retrieval-grounded paradigm
(\cite{mace,gautam,ted,veracity,bangare}) is more appropriate than the
classification paradigm~\cite{hemtfake} for three reasons. First, the system must
produce auditable, evidence-backed verdicts rather than opaque labels; the
classification approach optimises a decision boundary but cannot point to the
source that justifies a verdict. Second, the targeted misinformation is
fast-moving and local, so the system must reason over freshly retrieved evidence
rather than over patterns frozen at training time --- the explicit goal of the
retrieval-grounded systems. Third, the modular, agent-based designs allow each
stage to be evaluated and improved independently, which suits a production
pipeline. Among the agentic works, the verify-and-revise loop of MACE~\cite{mace}
and the provenance emphasis of Gautam~\cite{gautam} are the most directly
transferable, while Veracity~\cite{veracity} validates the end-user,
transparency-first framing. HEMT-Fake~\cite{hemtfake} remains relevant as a
multilingual and explainability benchmark and as a contrasting design point.

\subsection{Comparison with the Proposed Architecture}
Mapping the proposed system onto the literature reveals substantial alignment with
established practice alongside several distinguishing choices.

\begin{itemize}
    \item \textbf{Multi-agent orchestration.} The LangGraph pipeline --- a
    Claim-Extractor, a routing stage, Local- and Web-Researcher agents, and a
    Synthesis/Adjudication agent --- mirrors the decomposition logic of
    MACE~\cite{mace} and Gautam~\cite{gautam}. The Synthesis-and-Adjudication stage
    with a consistency cross-check is functionally the analogue of the MACE
    Verifier and the TED Synthesis agent~\cite{ted}.
    \item \textbf{Evidence grounding.} The Local- and Web-Researcher split
    parallels the web-retrieval agents of Veracity~\cite{veracity} and the
    Extractor of Gautam~\cite{gautam}, but the proposed system adds a hybrid local
    index (ChromaDB dense retrieval plus BM25 lexical matching) consulted
    \textit{before} the web.
    \item \textbf{External-API verification.} The use of an LLM for reasoning plus
    authoritative external sources echoes the Gemini-plus-fact-check-API pattern of
    Bangare et al.~\cite{bangare}, here substituting a domain-filtered web search
    over trusted Sri Lankan news and government domains.
    \item \textbf{Multilingual output.} The trilingual EN/SI/TA verdict aligns with
    the multilingual ambitions of Vykopal et al.~\cite{vykopal} and
    HEMT-Fake~\cite{hemtfake}, and the translate-toward-English insight of
    Vykopal et al. supports the decision to normalise code-mixed input into English
    claims.
    \item \textbf{Claim extraction.} The fine-tuned Phi-3.5 GGUF extractor, which
    converts noisy code-mixed posts into normalised declarative English claims, has
    the weakest precedent. FACT-GPT~\cite{factgpt} fine-tunes for claim
    \textit{matching}, not extraction or normalisation, and treats claims as
    already-formed inputs.
\end{itemize}

\subsection{Research Gaps}
Three gaps emerge, modest but real.

\begin{itemize}
    \item \textbf{Gap 1 --- Claim extraction from code-mixed, low-resource input as
    a dedicated stage.} The literature treats the upstream problem as claim
    \textit{matching}~\cite{factgpt} or claim \textit{detection}~\cite{vykopal},
    assuming a reasonably well-formed input claim. None addresses the normalisation
    of noisy, code-mixed Singlish/Sinhala/Tamil/English text into clean declarative
    claims as a distinct, fine-tuned step. Vykopal et al.~\cite{vykopal} show
    translation toward English helps, but build no extraction model;
    HEMT-Fake~\cite{hemtfake} is multilingual but classification-only and targets
    Indian languages.
    \item \textbf{Gap 2 --- Local-first retrieval with continuous OSINT ingestion
    and score-based escalation.} Veracity~\cite{veracity} and
    Bangare et al.~\cite{bangare} route claims directly to the open web or to a
    fact-check API. The proposed system instead maintains a continuously ingested
    local knowledge base (Telegram OSINT vectorised into ChromaDB with a
    synchronised BM25 index) and escalates to the web only when local evidence is
    insufficient. This local-first, score-gated routing is largely unexplored in
    the surveyed agentic systems.
    \item \textbf{Gap 3 --- Trilingual verdict generation for genuinely
    low-resource languages.} The agentic camp reasons well but is largely
    English-centric; the multilingual camp~\cite{vykopal,hemtfake} covers more
    languages but performs detection or matching, not evidence-grounded verdict
    generation with explanations. Producing adjudicated, explained verdicts in
    Sinhala and Tamil sits at an intersection the literature keeps separate.
\end{itemize}

\subsection{Contribution and Positioning}
This project occupies the underserved intersection of these gaps. It adopts the
well-validated agentic, retrieval-grounded paradigm --- modular agents, a
verify-and-revise adjudication step, and provenance-aware audit trails --- and
extends it in three directions the literature has not combined: a dedicated
fine-tuned claim-extraction stage for code-mixed low-resource input; a local-first
retrieval strategy backed by continuous OSINT ingestion with score-based web
escalation; and trilingual, evidence-grounded verdict generation for Sinhala and
Tamil. None of these is individually revolutionary, and the project deliberately
builds on established components rather than proposing a new agent topology; the
contribution is the integration of a low-resource, code-mixed claim pipeline into
a transparent multi-agent verification system targeted at a regional
misinformation ecosystem that prior work does not serve.
%=======================================================================================%

\section{Methodology}
\label{sec:methodology}

This section describes the end-to-end design and AI techniques employed in
TruthBomb. The system follows a linear multi-agent pipeline orchestrated by
LangGraph (a graph-based stateful agent framework built on LangChain), in which a shared typed state object
(\texttt{AgentState}) flows through four sequential stages: claim extraction,
evidence retrieval, adjudication, and trilingual verdict generation. Each
stage is implemented as an independent agent node, with a deterministic
router controlling branching between local and web retrieval paths. The
overall pipeline is illustrated conceptually in Figure~\ref{fig:pipeline}.

\begin{figure}[H]
\centering
\fbox{\parbox{0.85\textwidth}{%
\centering
\textbf{Input (code-mixed post)}\\[4pt]
$\downarrow$\\[2pt]
\textbf{Stage 1: Claim Extractor} \quad (Gemini 2.5 Flash Lite)\\
language ID $\cdot$ claim normalisation $\cdot$ routing decision\\[4pt]
$\downarrow$\\[2pt]
\textbf{Router} \quad (deterministic, confidence-based)\\[4pt]
\begin{tabular}{ccc}
  $\swarrow$ & & $\searrow$ \\
  \textbf{Stage 2a} & & \textbf{Stage 2b} \\
  Local Researcher & & Web Researcher \\
  (ChromaDB + Phi-3.5-mini) & & (DuckDuckGo + Gemini) \\
\end{tabular}\\[4pt]
$\searrow \quad \swarrow$\\[2pt]
\textbf{Stage 3: Synthesis Agent} \quad (Gemini 2.5 Flash Lite)\\
verdict $\cdot$ confidence $\cdot$ rationale $\cdot$ trilingual summary\\[4pt]
$\downarrow$\\[2pt]
\textbf{Output (EN / SI / TA)}
}}
\caption{High-level architecture of the TruthBomb multi-agent pipeline.}
\label{fig:pipeline}
\end{figure}

%-------------------------------------------------------------------
\subsection{Stage 1 -- Claim Extraction and Language Normalization}
\label{subsec:extraction}

The first node in the LangGraph workflow is the Claim Extractor agent. Its
responsibility is threefold: (i)~identify the primary language of the input
post, (ii)~isolate and normalize the single verifiable core claim into a
clean, professional English declarative sentence, and (iii)~emit a routing
signal that determines which retrieval branch the pipeline will follow.

The agent is powered by \textbf{Gemini 2.5 Flash Lite}, Google's lightweight
instruction-following model accessed via the \texttt{google.generativeai}
SDK. A structured prompt instructs the model to strip emotional language,
social-media noise, and opinions from the raw post, retaining only the
factual assertion that can be checked against evidence. The model is
constrained to return output exclusively as a JSON object with three fields:
\texttt{language} (one of \textit{English}, \textit{Sinhala},
\textit{Tamil}, or \textit{Singlish/Tanglish}), \texttt{extracted\_claim}
(the normalized English claim), and \texttt{routing\_destination} (either
\texttt{local\_researcher} or \texttt{web\_researcher}).

The rationale for framing extraction as normalization rather than translation
is consistent with the findings of Vykopal et al.~\cite{vykopal}: a targeted
intermediate representation toward English demonstrably improves downstream
reasoning quality on low-resource and code-mixed inputs. Critically, this
normalization is performed in a single inference step rather than two
separate translation and extraction passes, which avoids compounding errors
and reduces latency.

The routing decision embedded in this output is also deliberate. Claims about
local Sri Lankan events, government announcements, or incidents that are
likely represented in the continuously ingested OSINT database are directed
to the Local Researcher. Claims about international events or topics
requiring broad real-time internet context are routed to the Web Researcher.
This design mirrors the routing semantics of MACE~\cite{mace}, in which a
Planner decides strategy before execution, here collapsed into the extractor
itself to eliminate an additional inference step.

%-------------------------------------------------------------------
\subsection{Stage 2a -- Local Evidence Retrieval via Hybrid Search}
\label{subsec:local-retrieval}

When the router directs execution to the Local Researcher branch, the agent
queries a continuously maintained local OSINT knowledge base using a
\textbf{hybrid retrieval strategy} combining dense vector search and lexical
matching.

\paragraph{Knowledge base construction.}
The knowledge base is built from live Telegram channel messages ingested via
the Telethon MTProto client, polling three high-credibility Sri Lankan news
channels (\texttt{AdaDeranaNews}, \texttt{newsfirstsl},
\texttt{hirunewsofficial}) at 60-second intervals as a background
asynchronous task in FastAPI. Each ingested message is stored as a LangChain
\texttt{Document} object with structured metadata (source channel, message
ID, timestamp) and embedded using the
\textbf{paraphrase-multilingual-MiniLM-L12-v2} sentence-transformer model
from the \texttt{sentence-transformers} library. This model was chosen
because it produces semantically meaningful embeddings for multilingual text,
including Sinhala and Tamil, within a 384-dimensional dense vector space.
Embedded documents are persisted in a \textbf{ChromaDB} vector store.

A \textbf{BM25 lexical index} (via LangChain's \texttt{BM25Retriever}) is
maintained in parallel and updated every time new documents are added to
ChromaDB. This dual-index architecture implements the hybrid retrieval
pattern: dense retrieval captures semantic similarity even when surface
vocabulary differs, while BM25 captures keyword overlap that dense embeddings
can miss for short, highly specific named-entity queries (road names, place
names, official designations).

\paragraph{Retrieval and reasoning.}
At query time, the normalized claim is submitted to ChromaDB with $k{=}3$
for dense similarity search. The top-3 retrieved documents are concatenated
into a context block and passed, together with the claim, to a
\textbf{fine-tuned Phi-3.5-mini-instruct} model (Microsoft, 2024). The base model
(\texttt{unsloth/Phi-3.5-mini-instruct}) is loaded in \textbf{4-bit
quantization} using the BitsAndBytes library with \texttt{torch.float16}
precision to minimize GPU memory footprint, then augmented with a
\textbf{LoRA} (Low-Rank Adaptation) adapter
(\texttt{models/phi3-extractor-lora}) loaded via PEFT. The LoRA adapter was
fine-tuned on Sri Lanka-specific claim-verification data to specialize the
model's reasoning for the domain. The model receives a prompt in the
Phi-3.5-mini chat template format and generates a verification
summary of up to 512 tokens, which is returned as the Local Researcher's
findings to the shared state.

In the standalone FastAPI deployment (\texttt{src/main.py}), the same
knowledge base is queried by combining ChromaDB dense retrieval ($k{=}3$)
with BM25 ($k{=}2$), de-duplicated by document content, and passed to the
GGUF-quantized variant of the fine-tuned model (\texttt{phi3-extractor-unsloth.Q4\_K\_M.gguf})
running via \textbf{llama.cpp}. Structured output is enforced via a
\textbf{GBNF grammar} that constrains the model to emit only valid JSON
matching the schema \texttt{\{verdict, confidence\_score, reasoning\}},
eliminating post-hoc JSON repair and guaranteeing downstream parseability.

%-------------------------------------------------------------------
\subsection{Stage 2b -- Web Evidence Retrieval via Trusted-Domain Search}
\label{subsec:web-retrieval}

When the router directs execution to the Web Researcher branch, the agent
retrieves evidence from the open web while enforcing a domain restriction
policy to maintain source credibility.

The retrieval tool (implemented as a LangChain \texttt{@tool}) constructs a
DuckDuckGo query augmented with a compound \texttt{site:} filter covering
ten pre-vetted Sri Lankan and regional news domains (Ada Derana, NewsFirst,
Daily Mirror, Financial Times, Lanka Deepa, Virakesari, Economy Next, The
Island, \texttt{news.lk}, \texttt{government24.gov.lk}). This restricts
organic search results to authoritative outlets rather than the open web,
directly addressing the source credibility requirements identified by
Curtis et al.~\cite{veracity}. A fallback path executes an unrestricted
DuckDuckGo search if the domain-filtered query returns no results, and the
fallback status is propagated to the downstream synthesis agent so it may
weight that evidence with lower confidence.

The raw search results are then passed to \textbf{Gemini 2.5 Flash Lite}
with a structured prompt that instructs the model to: (i)~summarize the
findings from the retrieved articles, (ii)~explicitly attribute each
finding to its source domain, and (iii)~issue a preliminary verdict of
\textsc{supported}, \textsc{refuted}, or \textsc{unverified} based on the
retrieved content alone. This preliminary verdict is not exposed as a final
output; it is appended to the \texttt{agent\_outputs} list in the shared
state for the Synthesis agent to consume.

%-------------------------------------------------------------------
\subsection{Stage 3 -- Adjudication and Final Verdict Synthesis}
\label{subsec:synthesis}

After either retrieval branch completes, both paths converge at the
\textbf{Synthesis agent}, which is the terminal reasoning node in the
LangGraph workflow before \texttt{END}.

The Synthesis agent receives the normalized English claim and the full
\texttt{agent\_outputs} list from shared state, which contains the findings
and source metadata of whichever researcher executed. A single structured
prompt to \textbf{Gemini 2.5 Flash Lite} tasks the model with:

\begin{enumerate}
  \item Determining the final verdict from the four possible values:
        \textit{True}, \textit{False}, \textit{Misleading}, or
        \textit{Unverified}.
  \item Assigning a numerical confidence score (0--100\%).
  \item Producing a detailed English rationale that cites the evidence used
        to reach the verdict, providing an audit trail consistent with the
        provenance requirements articulated by Gautam~\cite{gautam}.
  \item Rendering a concise two-sentence verdict summary in each of three
        languages: \textbf{English}, \textbf{Sinhala} (සිංහල), and
        \textbf{Tamil} (தமிழ்).
\end{enumerate}

The model is constrained to return output as a JSON object with the schema
\texttt{\{verdict, confidence\_score, rationale, summaries:\{english,
sinhala, tamil\}\}}, which the agent parses and stores in
\texttt{final\_verdict} of the shared state. This structured output is then
consumed by the API layer and surfaced to the user interface.

The adjudication design is intentionally \emph{single-pass} rather than
iterative. Given the modular LangGraph architecture, a verify-and-revise
loop analogous to the MACE Verifier~\cite{mace} could be added as an
additional node, which remains a direction for future extension.

%-------------------------------------------------------------------
\subsection{Stage 4 -- Trilingual Verdict Delivery}
\label{subsec:trilingual}

The trilingual summaries produced by the Synthesis agent address a first-class
design requirement: the communities most exposed to Sri Lankan misinformation
are those who read in Sinhala and Tamil, not English. By generating the
verdict explanation natively in all three languages within a single inference
call, TruthBomb avoids the translation pipeline that Vykopal et
al.~\cite{vykopal} show degrades quality for low-resource languages. The
Gemini model's native multilingual generation capability, which includes
Sinhala and Tamil, makes this feasible without a separate translation agent.

Both frontend interfaces -- the React/Vite single-page application and the
Streamlit application -- present the English summary and rationale by
default, while the Sinhala and Tamil summaries are available in expandable
detail panels, making the output accessible to all target language communities.

%-------------------------------------------------------------------
\subsection{Orchestration and State Management}
\label{subsec:orchestration}

The multi-agent pipeline is orchestrated using \textbf{LangGraph}, a
graph-based extension of LangChain for building stateful, multi-step agent
workflows. The pipeline is defined as a \texttt{StateGraph} over the typed
\texttt{AgentState} dictionary, which carries all inter-agent communication:
\texttt{raw\_input}, \texttt{language}, \texttt{extracted\_claim},
\texttt{routing\_destination}, \texttt{agent\_outputs}, and
\texttt{final\_verdict}. Because state is passed by value rather than by
side effects, each agent node is a pure function of its state slice, which
simplifies debugging and enables deterministic re-execution.

Conditional routing is implemented via LangGraph's
\texttt{add\_conditional\_edges} primitive, where the router function reads
\texttt{routing\_destination} from state and returns the name of the next
node. Both the \texttt{local\_researcher} and \texttt{web\_researcher} nodes
connect unconditionally to \texttt{synthesis} via direct edges, ensuring the
pipeline always converges regardless of which retrieval branch executed.

The entire pipeline is exposed as a \textbf{FastAPI} REST endpoint
(\texttt{POST /api/v1/verify}), accepting a JSON payload
\texttt{\{text: str\}} and returning a structured \texttt{VerificationResponse}
that includes the raw input, extracted claim, verdict, confidence score,
retrieved context documents, reasoning, and timestamp. A health check
endpoint (\texttt{GET /api/v1/health}) reports database readiness, enabling
deployment monitoring.

%-------------------------------------------------------------------
\subsection{Summary of AI Techniques}
\label{subsec:ai-summary}

Table~\ref{tab:techniques} summarises the AI techniques employed at each
stage and the design rationale for each choice.

\begin{table}[H]
\centering
\caption{AI techniques used per pipeline stage.}
\label{tab:techniques}
\renewcommand{\arraystretch}{1.3}
\begin{tabular}{|p{3.5cm}|p{4.5cm}|p{5.5cm}|}
\hline
\textbf{Stage} & \textbf{Technique / Model} & \textbf{Rationale} \\
\hline
Claim extraction \& routing &
  Gemini 2.5 Flash Lite (zero-shot structured prompting) &
  Single-step normalization + routing; avoids cascading translation errors on
  code-mixed input~\cite{vykopal} \\
\hline
Local OSINT retrieval &
  Dense retrieval (ChromaDB + \texttt{paraphrase-multilingual-MiniLM-L12-v2}) +
  BM25 lexical retrieval &
  Hybrid search captures both semantic and keyword similarity; multilingual
  embeddings handle Sinhala and Tamil tokens \\
\hline
Local reasoning &
  Fine-tuned Phi-3.5-mini-instruct with LoRA (4-bit quantized via PEFT /
  GGUF via llama.cpp); GBNF grammar for structured output &
  Domain-adapted small model reduces hallucination; 4-bit quantization
  enables CPU/low-VRAM deployment; GBNF enforces JSON schema
  deterministically \\
\hline
Web retrieval &
  DuckDuckGo with trusted-domain \texttt{site:} filter &
  Restricts retrieved evidence to pre-vetted authoritative Sri Lankan sources;
  fallback propagated for downstream down-weighting \\
\hline
Web summarization &
  Gemini 2.5 Flash Lite (structured prompt) &
  Efficient, source-attributed summarization over variable-length search
  snippets \\
\hline
Adjudication \& synthesis &
  Gemini 2.5 Flash Lite (structured prompt + JSON output) &
  Combines evidence from any retrieval branch; produces auditable rationale
  consistent with provenance requirements~\cite{gautam} \\
\hline
Trilingual verdict &
  Gemini 2.5 Flash Lite (native multilingual generation) &
  Avoids translation pipeline degradation; native SI/TA generation in single
  pass~\cite{vykopal} \\
\hline
Continuous ingestion &
  Telethon MTProto + asynchronous FastAPI task &
  Low-latency Telegram polling keeps local knowledge base current without
  blocking the verification API \\
\hline
\end{tabular}
\end{table}

\section{Implementation}
\section{Experimental Setup}
\section{Performance Metrices}
\section{Evaluation}
\section{Discussion}
\section{Conclusions}

\bibliographystyle{IEEEtran}
\addcontentsline{toc}{section}{References}
\bibliography{references}

\end{document}